\documentclass[11pt,reqno]{amsart}

\usepackage{amsmath,amssymb,amsthm}
\usepackage[T1]{fontenc}
\usepackage[expansion=false]{microtype}
\usepackage[margin=1in]{geometry}
\usepackage{hyperref}
\hypersetup{colorlinks=true,linkcolor=blue,citecolor=blue,urlcolor=blue}

\theoremstyle{plain}
\newtheorem{theorem}{Theorem}[section]
\newtheorem{proposition}[theorem]{Proposition}
\newtheorem{lemma}[theorem]{Lemma}
\newtheorem{corollary}[theorem]{Corollary}

\theoremstyle{definition}
\newtheorem{definition}[theorem]{Definition}
\newtheorem{remark}[theorem]{Remark}

\newcommand{\Q}{\mathbb{Q}}
\newcommand{\Z}{\mathbb{Z}}
\DeclareMathOperator{\Jac}{Jac}

\sloppy

\begin{document}

\title[No perfect cuboid in the Saunderson family]{No perfect cuboid in the Saunderson family of Euler bricks}

\author{Lightman Chang}
\address{Independent Researcher}
\email{lightman.chang@gmail.com}

\subjclass[2020]{Primary 11D09; Secondary 11G05, 11G30}
\keywords{perfect cuboid, Euler brick, Saunderson family, elliptic curve, rank zero, Chabauty}

\date{\today}

\begin{abstract}
A perfect cuboid is a rectangular box whose three edges, three face diagonals,
and space diagonal are all integers; whether one exists is a question of Euler
(1769) that remains open. Saunderson, in the eighteenth century, exhibited a
two-parameter family of Euler bricks
$(a,b,c)=\bigl(u(4v^2-w^2),\,v(4u^2-w^2),\,4uvw\bigr)$ attached to a Pythagorean
triple $(u,v,w)$; this family produces infinitely many bricks but has never
yielded a perfect one. We prove that it never can. Substituting the standard
Pythagorean parametrization $(u,v,w)=(p^2-q^2,2pq,p^2+q^2)$ into the
square-space-diagonal condition produces, after an explicit algebraic
reduction, a genus-three curve; a palindromic substitution $W=t+1/t$ followed
by the lifting identity $W^2-4=(t-1/t)^2$ collapses the perfect-cuboid locus
onto the genus-one curve $C_0:S^2=T_0^4+72T_0^2+16$. The Jacobian of $C_0$ is the
elliptic curve $y^2=x^3-7x+6$, Cremona label \texttt{80a1}, of conductor $80$,
torsion $(\Z/2\Z)^2$, and Mordell--Weil rank zero (unconditionally, by
$2$-descent, and independently by Kolyvagin). Consequently $C_0(\Q)$ is finite
and consists of four points, all lying over $T_0\in\{0,\infty\}$, which we show
correspond to degenerate bricks with a vanishing edge. Hence no Saunderson
brick is a perfect cuboid. The result closes a single, explicitly delimited
subfamily of Euler bricks; it is complementary to the existence theorems of
Yoshida for face cuboids and is not covered by Peschmann's genus-three
quartic-pair reduction. All arithmetic claims are verified in PARI/GP.
\end{abstract}

\maketitle

\section{Introduction}\label{sec:intro}

A \emph{perfect cuboid} is a rectangular parallelepiped with positive integer
edges $a,b,c$ whose three face diagonals $\sqrt{a^2+b^2}$, $\sqrt{b^2+c^2}$,
$\sqrt{a^2+c^2}$ and whose space diagonal $\sqrt{a^2+b^2+c^2}$ are all integers.
The Perfect Cuboid Problem (PCP), which asks whether a perfect cuboid exists,
was raised by Euler in 1769 and is still open after more than two and a half
centuries. The strongest computational evidence comes from the absence of any
perfect cuboid with smallest edge below $5\times10^{11}$, or with odd edge below
$2.5\times10^{13}$~\cite{matson}.

An \emph{Euler brick} is the weaker object in which only the three face
diagonals are required to be integers; the perfect-cuboid problem is then the
question of whether some Euler brick also has an integral space diagonal.
Euler bricks are plentiful---the smallest, $(240,252,275)$, is due to Halcke
(1719)---but no parametrization of \emph{all} of them is known, and this is one
of the principal obstacles to a complete resolution of PCP. A natural strategy
is therefore to dispatch explicitly parametrized subfamilies one at a time.
Saunderson, in the eighteenth century, attached to each Pythagorean triple
$(u,v,w)$ the Euler brick
\begin{equation}\label{eq:saund}
\bigl(a,b,c\bigr)=\bigl(u(4v^2-w^2),\;v(4u^2-w^2),\;4uvw\bigr),
\qquad u^2+v^2=w^2 .
\end{equation}
Two threads in the recent literature bear directly on the elliptic-curve
geometry of such families. Yoshida~\cite{yoshida} studies \emph{face cuboids}
(only one face diagonal allowed to be irrational) through the elliptic family
$E_{1,s}:y^2=x(x-(2s)^2)(x+(s^2-1)^2)$, classifies its rational torsion as
$\Z/2\Z\times\Z/4\Z$, and uses non-torsion points to \emph{produce} infinitely
many face cuboids---an existence direction opposite to ours.
Peschmann~\cite{peschmann} reduces the full perfect-cuboid problem to a
one-parameter
family of genus-three curves $C_A:w^2=\lambda^8+A\lambda^4+1$ and develops
$2$-descent and Kummer-character obstructions on a distinguished elliptic
quotient; he does not treat the Saunderson subfamily, nor the specific
rank-zero curve that controls it, and explicitly leaves open
(\cite[\S8]{peschmann}) the adaptation of Asiryan's rank-zero
$\Q(\sqrt2)$ method ``from irreducibility to simultaneous representability.''

In this paper we prove that the Saunderson family~\eqref{eq:saund} contains no
perfect cuboid. The proof is an explicit reduction: the square-space-diagonal
condition for a Saunderson brick is equivalent to a rational point on a
genus-three curve, which a palindromic substitution and the elementary lifting
identity $W^2-4=(t-1/t)^2$ collapse onto a single genus-one curve
$C_0:S^2=T_0^4+72T_0^2+16$. The Jacobian of $C_0$ is the elliptic curve
$y^2=x^3-7x+6$ (Cremona label \texttt{80a1}), whose Mordell--Weil rank over
$\Q$ is zero; hence $C_0(\Q)$ is finite, equal to four points, all of which we
identify with degenerate (zero-edge) bricks. The closure thus rests on a
forty-year-old entry in Cremona's tables rather than on any conjecture.

The paper is organized as follows. Section~\ref{sec:family} records the
Saunderson family, derives the body-diagonal identity, and carries out the
algebraic reduction to the genus-one curve $C_0$ together with the lifting
condition. Section~\ref{sec:C0} establishes the arithmetic of $C_0$: its
Jacobian is \texttt{80a1}, of rank zero, so $C_0(\Q)$ has exactly four points.
Section~\ref{sec:main} assembles these into the proof of the main theorem and
checks the degeneracy bookkeeping. Section~\ref{sec:remarks} delimits the scope
of the result---the Saunderson family accounts for only a small fraction of
known Euler bricks---and positions the work relative to Yoshida and Peschmann.
All polynomial identities and curve computations are verified in PARI/GP; the
scripts and their captured output accompany the paper.

\section{The Saunderson family and its reduction}\label{sec:family}

\subsection{The family and the body-diagonal identity}

\begin{definition}\label{def:saund}
Let $(u,v,w)$ be a primitive Pythagorean triple, $u^2+v^2=w^2$ with
$\gcd(u,v,w)=1$. The associated \emph{Saunderson brick} is the triple
\[
\mathrm{Sa}(u,v,w):=\bigl(u(4v^2-w^2),\;v(4u^2-w^2),\;4uvw\bigr).
\]
A Saunderson brick is \emph{perfect} if it has all edges nonzero and
$a^2+b^2+c^2$ is a perfect square, where $(a,b,c)=\mathrm{Sa}(u,v,w)$.
\end{definition}

That $\mathrm{Sa}(u,v,w)$ is an Euler brick---i.e.\ that $a^2+b^2$, $b^2+c^2$,
$a^2+c^2$ are squares whenever $u^2+v^2=w^2$---is classical and is not needed
below; we use only the body-diagonal identity, which we record next.

\begin{lemma}\label{lem:body}
With $(a,b,c)=\mathrm{Sa}(u,v,w)$ and $u^2+v^2=w^2$,
\[
a^2+b^2+c^2 \;=\; w^2\bigl(w^4+16\,u^2v^2\bigr).
\]
\end{lemma}

\begin{proof}
Substitute the universal Pythagorean parametrization
$(u,v,w)=(p^2-q^2,\,2pq,\,p^2+q^2)$, under which both sides become polynomials
in $\Z[p,q]$. A direct polynomial expansion then shows that their difference
$\bigl(a^2+b^2+c^2\bigr)-w^2\bigl(w^4+16u^2v^2\bigr)$ is identically zero in
$\Z[p,q]$ (script~\texttt{01\_algebraic\_identities.gp}, Step~1). Every
primitive Pythagorean triple arises from such a pair $(p,q)$ up to sign, and
both sides are polynomial in $(u,v,w)$ subject only to $u^2+v^2=w^2$; hence the
identity holds for all Pythagorean $(u,v,w)$.
\end{proof}

\begin{corollary}\label{cor:gw}
$\mathrm{Sa}(u,v,w)$ has integral space diagonal $g$ if and only if
$(g/w)^2=w^4+16u^2v^2$.
\end{corollary}

\begin{proof}
Immediate from Lemma~\ref{lem:body}: $a^2+b^2+c^2=g^2$ if and only if
$g^2=w^2(w^4+16u^2v^2)$, i.e.\ $(g/w)^2=w^4+16u^2v^2$.
\end{proof}

\subsection{Reduction to a genus-one curve}

We now substitute $(u,v,w)=(p^2-q^2,2pq,p^2+q^2)$ and set $t=p/q\in\Q$. The
following three identities, each a polynomial (or Laurent-polynomial) identity
verified symbolically in script~\texttt{01\_algebraic\_identities.gp} and
\texttt{05\_scaling\_check.gp}, drive the reduction.

\begin{lemma}\label{lem:octic}
With $(u,v,w)=(p^2-q^2,2pq,p^2+q^2)$ and $t=p/q$,
\[
w^4+16u^2v^2 \;=\; q^8\bigl(t^8+68t^6-122t^4+68t^2+1\bigr).
\]
Consequently, writing $T:=g/(w\,q^4)$, the condition of
Corollary~\ref{cor:gw} is equivalent to a rational point on the genus-three
curve
\[
C':\qquad T^2=t^8+68t^6-122t^4+68t^2+1 .
\]
\end{lemma}

\begin{proof}
The displayed identity is checked by expanding $w^4+16u^2v^2$ as a polynomial
in $p,q$, substituting $p=tq$, and dividing by $q^8$
(script~\texttt{01}, Step~2; script~\texttt{05}, first block). By
Corollary~\ref{cor:gw}, integrality of $g$ is equivalent to
$(g/w)^2=w^4+16u^2v^2=q^8\bigl(t^8+68t^6-122t^4+68t^2+1\bigr)$, i.e.\ to
$T^2=t^8+68t^6-122t^4+68t^2+1$ with $T=(g/w)/q^4=g/(wq^4)$. The octic on the
right has degree $8$ with distinct roots, so $C'$ has genus three.
\end{proof}

\begin{lemma}\label{lem:palindrome}
The octic is palindromic, and with $W=t+1/t$ one has the Laurent-polynomial
identity
\[
\frac{t^8+68t^6-122t^4+68t^2+1}{t^4}=W^4+64W^2-256 .
\]
Hence, writing $S:=T/t^2$, the curve $C'$ is equivalent to
\[
C'_{\mathrm{pal}}:\qquad S^2=W^4+64W^2-256 .
\]
\end{lemma}

\begin{proof}
Dividing the octic by $t^4$ gives
$t^4+t^{-4}+68(t^2+t^{-2})-122$, which equals $W^4+64W^2-256$ upon
substituting $W=t+1/t$ (so that $t^2+t^{-2}=W^2-2$ and
$t^4+t^{-4}=W^4-4W^2+2$); the difference is identically zero
(script~\texttt{01}, Step~3). Dividing $T^2=t^8+\cdots+1$ by $t^4$ gives
$(T/t^2)^2=W^4+64W^2-256$, i.e.\ $S^2=W^4+64W^2-256$ with $S=T/t^2$.
\end{proof}

The Jacobian of $C'_{\mathrm{pal}}$ is the elliptic curve
$E_{\mathrm{PCP}}:y^2=x^3+x^2-x+15$ of conductor $160$, $j$-invariant
$-64/25$, torsion $\Z/2\Z$, and Mordell--Weil rank $1$ with generator
$(-1,4)$; analytic rank $1$ together with Kolyvagin's theorem~\cite{kolyvagin} makes the
rank unconditional (script~\texttt{03\_EPCP\_and\_hyperell.gp}). Because
$E_{\mathrm{PCP}}$ has rank $1$, the curve $C'_{\mathrm{pal}}$ alone does not
force finiteness of the relevant rational points; the perfect-cuboid condition
supplies one further constraint, which lowers the rank to zero. This is the
content of the next lemma.

\begin{lemma}[Lifting condition]\label{lem:lift}
For $t\in\Q^\times$ and $W=t+1/t$ one has $W^2-4=(t-1/t)^2$. Conversely, if
$W,T_0\in\Q$ satisfy $W^2-4=T_0^2$, then $t:=(W+T_0)/2$ is rational and
satisfies $t+1/t=W$, $t-1/t=T_0$. In particular, a rational point
$(W,S)\in C'_{\mathrm{pal}}(\Q)$ arises from a rational $t$ \emph{if and only
if} $W^2-4$ is a rational square.
\end{lemma}

\begin{proof}
The identity $W^2-4=(t-1/t)^2$ is verified directly
(script~\texttt{01}, Step~3b). For the converse, if $W^2-4=T_0^2$ set
$t=(W+T_0)/2$; then $1/t=(W-T_0)/2$ because $t\cdot\tfrac{W-T_0}{2}
=\tfrac{(W+T_0)(W-T_0)}{4}=\tfrac{W^2-T_0^2}{4}=1$, whence $t+1/t=W$ and
$t-1/t=T_0$ (script~\texttt{04\_degeneracy\_bookkeeping.gp}, final block).
\end{proof}

\begin{proposition}\label{prop:C0}
A perfect Saunderson brick yields a rational point $(T_0,S)$ with $T_0\neq0$ on
the genus-one curve
\[
C_0:\qquad S^2=T_0^4+72T_0^2+16 ,
\]
where $T_0=t-1/t$ and $S=T/t^2$.
\end{proposition}

\begin{proof}
Let $(a,b,c)=\mathrm{Sa}(u,v,w)$ be a perfect Saunderson brick with
$(u,v,w)=(p^2-q^2,2pq,p^2+q^2)$ and $t=p/q$. By Corollary~\ref{cor:gw} and
Lemmas~\ref{lem:octic} and~\ref{lem:palindrome}, the integrality of the space
diagonal gives a rational point $(W,S)$ on $C'_{\mathrm{pal}}$ with $W=t+1/t$
and $S=T/t^2$. By Lemma~\ref{lem:lift}, $W^2-4=T_0^2$ with $T_0=t-1/t\in\Q$.
Substituting $W^2=T_0^2+4$ into $S^2=W^4+64W^2-256=(W^2)^2+64W^2-256$ gives
\[
S^2=(T_0^2+4)^2+64(T_0^2+4)-256=T_0^4+72T_0^2+16,
\]
an identity verified symbolically (script~\texttt{01}, Step~4). Thus
$(T_0,S)\in C_0(\Q)$. Finally, a perfect brick has all edges nonzero, so
$u,v\neq0$ and $u\neq\pm v$; hence $p\neq\pm q$ and $p,q\neq0$, so
$t\notin\{0,\pm1,\infty\}$ and $T_0=t-1/t\neq0$.
\end{proof}

\section{The genus-one curve \texorpdfstring{$C_0$}{C0} and its rank-zero Jacobian}\label{sec:C0}

\begin{proposition}\label{prop:jac}
The Jacobian of $C_0:S^2=T_0^4+72T_0^2+16$ is the elliptic curve
\[
E_0:\qquad y^2=x^3-7x+6
\]
of conductor $80$ \textup{(}Cremona label \texttt{80a1}\textup{)}, with
$j$-invariant $148176/25$ and torsion subgroup $(\Z/2\Z)^2$. The full
Mordell--Weil group is
\[
E_0(\Q)=\{\,O,\;(1,0),\;(2,0),\;(-3,0)\,\},
\]
and $E_0$ has Mordell--Weil rank $0$.
\end{proposition}

\begin{proof}
Applying the standard quartic-to-Weierstrass transformation to
$y^2=x^4+72x^2+16$ (which has the rational point $(0,4)$) yields, after
minimalization, the model $y^2=x^3-7x+6$; its conductor is $80$ and its
$j$-invariant is $148176/25$ (script~\texttt{02\_curve\_C0\_and\_80a1.gp}).
Since $x^3-7x+6=(x-1)(x-2)(x+3)$, the three nontrivial $2$-torsion points are
$(1,0),(2,0),(-3,0)$, so $(\Z/2\Z)^2\subseteq E_0(\Q)_{\mathrm{tors}}$; PARI's
\texttt{elltors} confirms the torsion is exactly $(\Z/2\Z)^2$.

The rank is computed by $2$-descent: PARI's \texttt{ellrank} returns the tight
interval $[0,0]$, a proof that $\mathrm{rank}\,E_0(\Q)=0$ (the algorithm is
$2$-descent, which is unconditional). Independently, $E_0$ has analytic rank
$0$ with $L(E_0,1)\approx1.0095\neq0$; by Kolyvagin's theorem this also forces
algebraic rank $0$. The two computations agree, and the conductor, torsion, and
rank match the Cremona-table entry \texttt{80a1}. With rank $0$ and torsion
$(\Z/2\Z)^2$ the group $E_0(\Q)$ has the four elements listed.
\end{proof}

\begin{remark}
As a cross-check on the identification, the genus-three curve $C'$ of
Lemma~\ref{lem:octic} carries the commuting involutions $\sigma:t\mapsto-t$ and
$\tau:t\mapsto1/t$, and its Jacobian decomposes up to isogeny as a product of
three elliptic quotients $X_\sigma\times X_\tau\times X_{\sigma\tau}$. One has
$X_\sigma\cong X_\tau\cong E_{\mathrm{PCP}}$ (conductor $160$, rank $1$), while
$X_{\sigma\tau}$ is exactly $E_0=\texttt{80a1}$: the Frobenius traces $a_p$ of
$E_0$ for $p=3,7,11,\dots,47$ are
\[
(0,4,-4,-2,2,-4,-4,-2,8,6,-6,8,-4),
\]
matching $a_p(C')-2a_p(E_{\mathrm{PCP}})$ in all $13$ cases
(script~\texttt{02}). Thus $\mathrm{rank}\,\Jac(C')=1+1+0=2<3=\mathrm{genus}(C')$,
so Chabauty's method applies to $C'$; the reduction of
Section~\ref{sec:family}, however, bypasses Chabauty by landing directly on the
genus-one curve $C_0$.
\end{remark}

\begin{proposition}\label{prop:C0pts}
The set $C_0(\Q)$ has exactly four elements:
\[
C_0(\Q)=\{\,(0,4),\;(0,-4),\;\infty_{+},\;\infty_{-}\,\},
\]
where $(0,\pm4)$ are the affine points with $T_0=0$ and $\infty_{\pm}$ are the
two points at infinity. In particular every rational point of $C_0$ has
$T_0\in\{0,\infty\}$.
\end{proposition}

\begin{proof}
$C_0$ has the rational point $(0,4)$, hence is a smooth genus-one curve with a
rational point and is therefore isomorphic over $\Q$ to its Jacobian $E_0$.
By Proposition~\ref{prop:jac}, $|C_0(\Q)|=|E_0(\Q)|=4$.

We exhibit the four points. The leading coefficient of $T_0^4+72T_0^2+16$ is
$1=1^2$, a square, so the two points at infinity $\infty_{\pm}$ (where
$S/T_0^2\to\pm1$) are rational; these account for $T_0=\infty$. For the affine
points, an exhaustive search over $T_0=n/d$ in lowest terms with $|n|\le2000$,
$d\le2000$ finds $T_0=0$ as the only value for which $T_0^4+72T_0^2+16$ is a
square, giving $(0,\pm4)$ (script~\texttt{02}). An independent run of PARI's
\texttt{hyperellratpoints} on $x^4+72x^2+16$ with search bound $10^4$ returns
exactly the affine points $[0,4]$ and $[0,-4]$
(script~\texttt{03}). These two affine points together with $\infty_{\pm}$ are
four rational points; since $|C_0(\Q)|=4$, they are all of $C_0(\Q)$. Every one
of them has $T_0\in\{0,\infty\}$.
\end{proof}

\section{Proof of the Main Theorem}\label{sec:main}

\begin{theorem}\label{thm:main}
No Saunderson brick is a perfect cuboid. That is, for every primitive
Pythagorean triple $(u,v,w)$ the Saunderson brick
\[
\mathrm{Sa}(u,v,w)=\bigl(u(4v^2-w^2),\;v(4u^2-w^2),\;4uvw\bigr)
\]
fails to be a perfect cuboid: either some edge vanishes, or $a^2+b^2+c^2$ is
not a perfect square.
\end{theorem}

\begin{proof}
Suppose, for contradiction, that $\mathrm{Sa}(u,v,w)$ is a perfect cuboid in
the sense of Definition~\ref{def:saund}: all edges are nonzero and
$a^2+b^2+c^2$ is a perfect square. By Proposition~\ref{prop:C0} this produces a
rational point $(T_0,S)\in C_0(\Q)$ with $T_0\neq0$ and $T_0$ finite (since
$t=p/q\notin\{0,\pm1,\infty\}$ forces $T_0=t-1/t$ to be a nonzero rational, in
particular not $\infty$). But by Proposition~\ref{prop:C0pts} every rational
point of $C_0$ has $T_0\in\{0,\infty\}$, so no rational point of $C_0$ has
$T_0$ a nonzero finite rational. This contradiction proves the theorem.
\end{proof}

We record the degeneracy bookkeeping that identifies the four points of
$C_0(\Q)$ with degenerate bricks; it confirms that the only way the chain can
terminate is through a vanishing edge, consistent with the theorem.

\begin{proposition}\label{prop:deg}
Each rational point of $C_0$ corresponds to a degenerate Saunderson brick:
\begin{itemize}
\item $T_0=0$ forces $t=\pm1$, i.e.\ $p=\pm q$, hence $u=p^2-q^2=0$ and the
edge $a=u(4v^2-w^2)=0$;
\item $T_0=\infty$ forces $t\in\{0,\infty\}$, i.e.\ $p=0$ or $q=0$, hence
$v=2pq=0$, so the edges $b=v(4u^2-w^2)$ and $c=4uvw$ both vanish.
\end{itemize}
\end{proposition}

\begin{proof}
If $T_0=t-1/t=0$ then $t^2=1$, so $t=\pm1$ and $p=\pm q$; then
$u=p^2-q^2=0$ and $a=u(4v^2-w^2)=0$. If $T_0=\infty$ then $W=t+1/t=\infty$,
which requires $t\in\{0,\infty\}$, i.e.\ $p=0$ or $q=0$; in either case
$v=2pq=0$, so $b=v(4u^2-w^2)=0$ and $c=4uvw=0$. Numerical instances
($p=q=1$ giving $(a,b,c)=(0,-8,0)$, and $p=1,q=0$ giving $(-1,0,0)$) appear in
script~\texttt{04}.
\end{proof}

\begin{remark}[A worked non-example]
The smallest Saunderson brick is $\mathrm{Sa}(3,4,5)=(117,44,240)$ (the brick
$(44,117,240)$ up to order), obtained from $p=2,q=1$, $t=2$. Here
$T_0=t-1/t=3/2\neq0$, and indeed $T_0^4+72T_0^2+16=2929/16$ is not a rational
square, so $(T_0,S)\notin C_0(\Q)$; correspondingly $a^2+b^2+c^2=73225$ is not
a perfect square (script~\texttt{05}). This single example already illustrates
the mechanism: a non-degenerate Saunderson value $t$ yields a nonzero finite
$T_0$, which can never lie on the four-point set $C_0(\Q)$.
\end{remark}

\section{Remarks on scope and related work}\label{sec:remarks}

\subsection{Scope: the Saunderson family is a small subfamily}\label{sec:scope}

Theorem~\ref{thm:main} closes the Saunderson family, not the full
Perfect Cuboid Problem, and this family~\eqref{eq:saund} accounts for only
a small fraction of all Euler bricks. Among the primitive Euler bricks with all
edges at most $1000$, exactly one---$(44,117,240)=\mathrm{Sa}(3,4,5)$---is of
Saunderson form; the other four, including Halcke's $(240,252,275)$, are not
(script~\texttt{03}). Extending the search to $5000$ leaves $9$ of $11$
primitive Euler bricks outside the Saunderson family. Thus a perfect cuboid
arising from a non-Saunderson Euler brick would be entirely invisible to the
reduction of Section~\ref{sec:family}, and the present result must not be read
as a resolution of PCP. Closing further families requires either a uniform
parametrization of all Euler bricks (not currently available) or a separate
arithmetic treatment of each family; the latter can in principle be carried
out family by family by analogous rank-zero arguments, but no uniform
statement is known.

\subsection{Relation to Yoshida and Peschmann}\label{sec:related}

Saunderson's parametrization~\eqref{eq:saund} is classical. The elliptic curve
$E_{\mathrm{PCP}}$ appearing after the palindromic step
(Lemma~\ref{lem:palindrome}) is, on the Pythagorean locus, isomorphic to the
curve family $E_{1,s}$ studied by Yoshida~\cite{yoshida}; in particular the
rank-$1$ structure and the torsion classification of that family are shared.
Yoshida's results are existence statements about \emph{face} cuboids and run in
the opposite direction from ours: he uses non-torsion points to produce
infinitely many face cuboids, whereas we use a rank-\emph{zero} curve to
exclude perfect cuboids on a parametric locus. The two are complementary.
What is specific to the present paper is the second, rank-lowering step: the
lifting condition $W^2-4=T_0^2$ (Lemma~\ref{lem:lift}) that collapses the
rank-$1$ curve $C'_{\mathrm{pal}}$ onto the rank-$0$ curve $C_0$. This step has
no counterpart in~\cite{yoshida}.

Peschmann~\cite{peschmann} reduces the full perfect-cuboid problem to a
genus-three quartic-pair family $C_A:w^2=\lambda^8+A\lambda^4+1$ and studies a
distinguished elliptic quotient via $2$-descent and a Kummer character. He does
not treat the Saunderson subfamily, and the specific rank-zero curve
\texttt{80a1} that controls it does not appear in his work; the only explicit
rank-zero curve in~\cite{peschmann} is Asiryan's
$Y^2=X(X-8)(X-9)$, invoked for an irreducibility argument, and the adaptation
of that approach ``from irreducibility to simultaneous representability'' is
listed there as open (\cite[\S8]{peschmann}). The Saunderson closure presented
here is therefore not covered by Peschmann's reduction.

\subsection{On novelty}\label{sec:novelty}

To the best of our knowledge, the explicit statement that the Saunderson family
contains no perfect cuboid, and in particular its reduction to the rank-zero
curve \texttt{80a1}, has not appeared in the literature. The torsion structure
of the related elliptic family is known
(\cite{yoshida}); the rank-zero curve \texttt{80a1} is a routine entry in
Cremona's tables~\cite{cremona}; and Saunderson's parametrization is classical.
What is new is the assembly: the body-diagonal identity of Lemma~\ref{lem:body}
specific to the Saunderson form, the palindromic reduction, and the lifting
step that lowers the rank from one to zero and thereby forces finiteness
without recourse to Chabauty or to any conjecture. We do not claim that the
underlying rank-zero phenomenon is unknown to specialists as folklore; we claim
only that the precise statement and its proof do not appear in the cited
literature.

\section*{Computational verification}

All polynomial identities (Lemmas~\ref{lem:body}, \ref{lem:octic},
\ref{lem:palindrome}, \ref{lem:lift}, and Proposition~\ref{prop:C0}) were
verified as exact identities in $\Z[p,q]$ or $\Q(t)$ in PARI/GP, and all curve
data (conductor, $j$-invariant, torsion, rank by both $2$-descent and analytic
rank, the $a_p$ match, and the rational points of $C_0$) were computed in
PARI/GP version~2.15.4~\cite{pari}. The scripts
\texttt{01\_algebraic\_identities.gp},
\texttt{02\_curve\_C0\_and\_80a1.gp},
\texttt{03\_EPCP\_and\_hyperell.gp},
\texttt{04\_degeneracy\_bookkeeping.gp}, and
\texttt{05\_scaling\_check.gp}, together with their captured \texttt{.out}
files, accompany this paper.

\section*{Acknowledgements}
The author thanks the maintainers of PARI/GP.

\begin{thebibliography}{9}

\bibitem{cremona}
J.~E. Cremona,
\emph{Algorithms for modular elliptic curves},
2nd ed., Cambridge University Press, Cambridge, 1997.
Online tables: curve \texttt{80a1} (Cremona) $=$ \texttt{80.a2} (LMFDB),
\url{https://www.lmfdb.org/EllipticCurve/Q/80/a/2}.

\bibitem{kolyvagin}
V.~A. Kolyvagin,
\emph{Euler systems},
in: The Grothendieck Festschrift, Vol.~II,
Progr. Math. \textbf{87}, Birkh\"auser, Boston, 1990, pp.~435--483.

\bibitem{matson}
R.~D. Matson,
\emph{Results of a computer search for a perfect cuboid},
unpublished manuscript, 2009;
\url{https://www.unsolvedproblems.org/}.

\bibitem{pari}
The PARI~Group,
\emph{PARI/GP version 2.15.4},
Univ. Bordeaux, 2023,
\url{https://pari.math.u-bordeaux.fr/}.

\bibitem{peschmann}
R.~Peschmann,
\emph{Quartic reductions and elliptic obstructions for perfect Euler bricks},
preprint, arXiv:2604.09328 (2026).

\bibitem{yoshida}
T.~Yoshida,
\emph{The relationship between face cuboids and elliptic curves},
preprint, arXiv:2407.09825 (2024).

\end{thebibliography}

\end{document}
